\begin{abstract}
    We propose a new method, dubbed One Class Signed Distance Function (OCSDF), to perform One Class Classification (OCC) by provably learning the Signed Distance Function (SDF) to the boundary of the support of any distribution. The distance to the support can be interpreted as a normality score, and its approximation using 1-Lipschitz neural networks provides robustness bounds against $l2$ adversarial attacks, an under-explored weakness of deep learning-based OCC algorithms. As a result, OCSDF comes with a new metric, certified AUROC, that can be computed at the same cost as any classical AUROC. We show that OCSDF is competitive against concurrent methods on tabular and image data while being way more robust to adversarial attacks, illustrating its theoretical properties. Finally, as exploratory research perspectives, we theoretically and empirically show how OCSDF connects OCC with image generation and implicit neural surface parametrization.
\end{abstract}


%\begin{abstract}
%We propose a new Novelty Detection and One Class classifier, based on the smoothness properties of orthogonal neural network, and on the properties of Hinge Kantorovich Rubinstein (HKR) function. The classifier benefits from robustness certificates against $l2$-attacks thanks to the Lipschitz constraint, whilst the HKR loss allows to provably approximate the signed distance function to the boundary of the distribution. The normality score induces by the classifier has a meaningful interpretation in term of distance to the support. Finally, gradient steps in the input space allows free generation of samples from the one class support in a fashion that reminds GAN or diffusion models. This opens path to certifiably robust and explainable One Class classifiers. The performance of our method is assessed on image datasets and tabular data, both in One Class and Anomaly Detection settings. The theoretical properties are corroborated by experiments. \footnote{Our code is available at {\small\url{anonymous.4open.science/r/todo}}}  
%\end{abstract}



    %We propose a new method that provably learns the Signed Distance Function (SDF) to the boundary of the support of any distribution, using 1-Lipschitz orthogonal neural classifier. The distance to the support can be interpreted as a normality score which allows to use the model for One Class Classification (OCC). On tabular data we consistently beat concurrent methods based on deep learning, and we match the results of tree based models. In 3D computer vision we recover smooth Neural Implicit Surfaces from 3D point clouds. The algorithm is able to leverage labelled data in OC and AD when available, and any human prior. This allows to deploy the model in industrial setting where the dataset of known anomalies increases with time. Finally, gradient steps in the input space allows free generation of samples from the support, in a fashion that reminds GAN or diffusion models. This opens a path to certifiably robust and explainable One Class classifiers. Our code is available at \url{https://anonymous.4open.science/r/CSDFL}